%! Symmetric Positive Definite Matrices — Unit 6, Lesson 6.3 — Student Packet Cover Sheet
\documentclass[10pt]{article}
\usepackage{linalg-article}
\usepackage{linalg-boxes}
\usepackage{ltablex}
\keepXColumns

\begin{document}

% Full-bleed burgundy banner
\begin{tikzpicture}[remember picture, overlay]
  \fill[burgundy] (current page.north west)
    rectangle ([yshift=-0.9in]current page.north east);
\end{tikzpicture}
\vspace{-0.8in}

\begin{center}
  {\color{white}\textbf{\LARGE Linear Algebra}}\\[4pt]
  {\color{blushmid}\large\textbf{Unit 6 \quad Eigenvalues and Eigenvectors}}\\[6pt]
  {\color{white}\textbf{\large Lesson 6.3 \quad Symmetric Positive Definite Matrices}}
\end{center}

\vspace{0.22in}
\namedateperiod
\vspace{0.18in}

\begin{learningtargetbox}
\small
\begin{enumerate}[leftmargin=1.5em, itemsep=5pt]
  \item \ldots{} recognize a \textbf{symmetric} matrix ($S=S^{\mathsf T}$) and explain why its eigenvectors are \emph{perpendicular}, so it diagonalizes as $S=Q\Lambda Q^{\mathsf T}$ with an orthonormal $Q$.
  \item \ldots{} test whether a symmetric matrix is \textbf{positive definite} --- all eigenvalues positive, or $a>0$ and $ac-b^2>0$.
  \item \ldots{} connect positive definiteness to positive \textbf{energy} $\mathbf{x}^{\mathsf T}S\mathbf{x}>0$, and explain why such a matrix is always invertible.
\end{enumerate}
\end{learningtargetbox}

\vspace{0.14in}

\begin{tocbox}
\small
\renewcommand{\arraystretch}{1.5}
\begin{tabularx}{\linewidth}{@{} c l X r @{}}
\rowcolor{blush}
\textbf{\#} & \textbf{Component} & \textbf{Description} & \textbf{Score} \\
\midrule
1 & Warm-Up        & Spiral review: the transpose \& $S=S^{\mathsf T}$, the dot-product $\perp$ test, normalizing to unit length & \blank{1.2cm} \\
2 & Guided Notes   & Symmetric $\Rightarrow\perp$ eigenvectors; $S=Q\Lambda Q^{\mathsf T}$; positive definite \& energy & \blank{1.2cm} \\
3 & Group Activity & Confirm symmetry \& $\perp$; a full $S=Q\Lambda Q^{\mathsf T}$; a symmetric matrix that is \emph{not} positive definite & \blank{1.2cm} \\
4 & Exit Ticket    & Quick check: symmetry, perpendicular eigenvectors, and the positive-definite test & \blank{1.2cm} \\
5 & Homework       & Spectral factorization, a positive-definite screen, energy, and justification & \blank{1.2cm} \\
\midrule
  & \multicolumn{2}{r}{\textbf{Total}} & \blank{1.2cm} \\
\end{tabularx}
\end{tocbox}

\vspace{0.10in}

\begin{remindbox}
\small In Lesson~6.2 a matrix diagonalized only when it had enough independent eigenvectors, and computing
$X^{-1}$ could be messy. \textbf{Symmetric} matrices ($S=S^{\mathsf T}$) fix both problems at once: their
eigenvectors always come out \emph{perpendicular}, so after normalizing them into an orthonormal $Q$ the
inverse is just the transpose --- $S=Q\Lambda Q^{\mathsf T}$. When every eigenvalue is also \emph{positive},
$S$ is \textbf{positive definite}: it stores positive energy $\mathbf{x}^{\mathsf T}S\mathbf{x}>0$ and is
guaranteed invertible.
\end{remindbox}

\end{document}
